\section{引言：Scaling Laws 的实践应用}

本节课是 Scaling Laws 系列的第二讲，也是最后一讲。与上一讲偏重理论基础不同，本讲更注重\textbf{案例研究}和\textbf{实践细节}。课程分为两个核心部分：

\begin{enumerate}
    \item \textbf{案例研究}：分析多个实际模型（Cerebras GPT、MiniCPM、DeepSeek LLM 等）如何在训练过程中应用 scaling laws 来做出设计决策
    \item \textbf{MuP（Maximal Update Parameterization）}：深入理解一种让超参数在不同模型规模间保持稳定的参数化方法
\end{enumerate}

\begin{figure}[H]
\centering
\includegraphics[width=0.95\textwidth]{slides-images/slide-000.jpg}
\caption{Lecture 11 封面：Scaling Laws 2\protect\footnotemark}
\end{figure}
\footnotetext{Slides 第1页。}

\begin{importantbox}{本讲的核心问题}
\begin{itemize}
    \item Chinchilla 的 scaling law 方法真的有效吗？在实际大模型训练中表现如何？
    \item 如何在小规模实验中确定学习率、batch size 等超参数，并可靠地迁移到大规模训练？
    \item 是否存在某种参数化方式，使得最优超参数不随模型规模变化？
\end{itemize}
\end{importantbox}

\begin{figure}[H]
\centering
\includegraphics[width=0.95\textwidth]{slides-images/slide-001.jpg}
\caption{课程动机回顾：scaling laws 用于指导大模型的超参数选择和架构设计\protect\footnotemark}
\end{figure}
\footnotetext{Slides 第2页。}

\begin{knowledgebox}{从 Chinchilla 到实践的知识鸿沟}
Chinchilla 论文之后，ChatGPT 的出现改变了大语言模型的竞争格局。各大前沿实验室不再公开发表关于 scaling 的详细研究。讲者提到，他与前沿实验室的人交流时，对方明确表示不会透露任何关于 scaling 的细节。因此，我们不得不依赖少数开放的模型训练报告来理解 scaling 在实践中的运作方式。
\end{knowledgebox}

\subsection{案例研究概览}

\begin{figure}[H]
\centering
\includegraphics[width=0.95\textwidth]{slides-images/slide-002.jpg}
\caption{案例研究覆盖的模型：Cerebras GPT、MiniCPM、DeepSeek LLM 三个核心案例，以及 Llama 3、Hunyuan Large、MiniMax-01 三个补充案例\protect\footnotemark}
\end{figure}
\footnotetext{Slides 第3页。}

讲者指出，在所有公开的 scaling 研究中，DeepSeek 和 MiniCPM 仍然是迄今为止\textbf{最详细的开放式 scaling 研究}——即使到了 2025 年，这一点也没有改变。

\subsection{本章小结}

Scaling laws 不仅是理论上的曲线拟合工具，更是指导实际大模型训练的核心方法论。本讲将通过多个真实案例，展示不同团队如何将 scaling laws 融入从超参数选择到模型规模决策的完整训练流程。

%% ========== Part 2: Cerebras GPT ==========

